% =====================================================================
\chapter{Superare l'esame: tutte le strade}
% =====================================================================

I moduli ti hanno insegnato il C++. Questo capitolo ti insegna l'\textbf{esame}: come è
fatto davvero, dove si perdono i punti, cosa fare minuto per minuto in ognuna delle tre
prove, quali aiuti puoi usare (e quasi nessuno usa) e cosa fare se una prova va male.
Molti studenti preparati vengono bocciati non perché non sanno il C++, ma perché non
sanno come funziona la prova: sprecano tempo sul quiz, consegnano un laboratorio che non
compila, all'orale non sanno spiegare il proprio codice. Sono tutti errori evitabili.

\begin{trappola}[: regole da verificare]
Le regole che seguono vengono dalla scheda ufficiale del corso e dalle slide del docente
(entrambe nel file \texttt{Prog2OMNIBUSpdf.pdf}). Le regole possono cambiare da un anno
all'altro: prima dell'appello controlla la \textbf{pagina avvisi del docente} e, se hai un
dubbio, chiedi a ricevimento. In questo capitolo le cose \q{da verificare} sono segnalate
esplicitamente.
\end{trappola}

% ---------------------------------------------------------------------
\section{Come funziona davvero l'esame}
% ---------------------------------------------------------------------

\begin{analogia}[: le tre porte]
Immagina un corridoio con tre porte in fila. Ogni porta si apre solo con la chiave che
trovi dietro la porta precedente. La prima porta (il \textbf{quiz}) ti dà la chiave per
entrare in laboratorio. Il \textbf{laboratorio} ti dà la chiave per l'orale, ma
\emph{solo} se il programma che consegni funziona. L'\textbf{orale} è la stanza dove si
decide il voto. Non serve essere perfetti alla prima porta: serve passarla. Il voto si
costruisce soprattutto nelle ultime due.
\end{analogia}

\begin{center}
\begin{tikzpicture}[node distance=1.9cm,
    porta/.style={processo, text width=2.9cm, minimum height=1.9cm}]
  \node[porta] (quiz) {\textbf{1. Quiz}\\30 domande\\1 ora};
  \node[porta, right=of quiz] (lab) {\textbf{2. Laboratorio}\\un programma C++\\circa 2 ore};
  \node[porta, right=of lab] (orale) {\textbf{3. Orale}\\il tuo codice\\+ teoria};
  \draw[freccia] (quiz) -- node[above, etichetta] {superato} (lab);
  \draw[freccia] (lab) -- node[above, etichetta] {funziona} (orale);
  \node[below=1.1cm of quiz, etichetta, text width=2.6cm, align=center] (no1)
       {non superato:\\ritenti all'appello successivo};
  \node[below=1.1cm of lab, etichetta, text width=2.6cm, align=center] (no2)
       {non funziona:\\niente orale};
  \node[below=1.1cm of orale, etichetta, text width=2.6cm, align=center] (voto)
       {voto finale:\\valutazione complessiva};
  \draw[frecciarossa] (quiz) -- (no1);
  \draw[frecciarossa] (lab) -- (no2);
  \draw[freccia] (orale) -- (voto);
\end{tikzpicture}
\end{center}

Ecco le regole, una prova alla volta, con le parole della scheda del corso e delle slide
del docente.

\begin{center}\small
\begin{tabular}{lp{5.2cm}p{6.2cm}}
\toprule
\textbf{Prova} & \textbf{Com'è fatta} & \textbf{Cosa significa per te} \\
\midrule
Quiz & 30 quesiti a risposta multipla su \emph{tutto} il programma, durata 1 ora. Tutti i
fogli usati vanno consegnati. &
Circa 2 minuti a domanda. È uno \textbf{sbarramento}: solo chi lo supera accede al
laboratorio. Puoi (e devi) usare i fogli per disegnare e simulare il codice. \\
\addlinespace
Laboratorio & Un problema da risolvere scrivendo un programma C++. Circa 2 ore (la scheda
dice \q{tipicamente 2/3 ore}). Si compila con \texttt{g++} su una macchina generica. Si può
consultare la documentazione del C++, \textbf{niente Internet}. &
Allenati \textbf{senza Internet e senza AI}, con il compilatore da terminale. Accede
all'orale solo chi consegna una soluzione \textbf{funzionante}: un programma che non
compila ti ferma qui. \\
\addlinespace
Orale & Si compila e si esegue il tuo progetto, poi domande sull'implementazione e su
tutto il programma. &
Devi saper spiegare \emph{ogni riga} che hai scritto e \emph{perché} l'hai scritta così.
Poi teoria di tutti i 14 moduli. \\
\bottomrule
\end{tabular}
\end{center}

Altre regole pratiche che vengono dalle stesse fonti:
\begin{itemize}
  \item Di solito le tre prove si fanno \textbf{nello stesso giorno}, se il numero di
        prenotati lo permette. Quindi preparati a una giornata lunga: porta acqua e
        qualcosa da mangiare.
  \item La \textbf{prenotazione su SmartEdu è obbligatoria}. Senza prenotazione non sostieni
        l'esame, anche se sei preparatissimo. Segna sul calendario l'apertura e la
        chiusura delle prenotazioni.
  \item Propedeuticità: serve aver superato Programmazione 1 (tu ce l'hai).
  \item Il voto finale è \q{una valutazione complessiva delle tre prove}: non è una somma
        aritmetica. Un laboratorio molto buono e un orale chiaro possono compensare un quiz
        appena sufficiente.
\end{itemize}

\begin{consiglio}[: tre cose da chiedere al docente (o ai tutor) prima dell'appello]
\begin{enumerate}
  \item Nel quiz le risposte \textbf{sbagliate tolgono punti}? Quante risposte giuste
        servono per passare? (La scheda parla di un minimo del 60\%: su 30 domande sono
        18.)
  \item In laboratorio, \textbf{quale documentazione} si può consultare e in che forma
        (pagine offline sul PC del laboratorio, un libro)?
  \item Come si \textbf{consegna} il laboratorio? Negli esami del 2022 la consegna avveniva
        con \texttt{git} su GitHub Classroom (clone, add, commit, push). Se è ancora così,
        devi saper usare quei quattro comandi senza pensarci.
\end{enumerate}
Ognuna di queste risposte cambia la tua strategia: è il miglior investimento di 10 minuti
che puoi fare.
\end{consiglio}

\begin{trappola}[: frequenza e studenti lavoratori]
La scheda del corso dice che la frequenza \q{è di norma obbligatoria}. Se lavori e non hai
frequentato, \textbf{verifica subito} con la segreteria didattica o con il docente se
questo ti riguarda e se esistono regole per gli studenti lavoratori. Meglio scoprirlo
oggi che il giorno dell'appello.
\end{trappola}

\begin{inparole}
Tre prove in fila, ognuna è un cancello per la successiva. Il quiz va solo
\emph{superato}. Il laboratorio deve \emph{funzionare}. All'orale si decide il voto,
partendo dal tuo codice. Prenotati su SmartEdu e chiedi al docente le tre regole che
mancano.
\end{inparole}

% ---------------------------------------------------------------------
\section{Prova 1 -- Il quiz: superare la prima porta}
% ---------------------------------------------------------------------

Il quiz ha 30 domande e un'ora di tempo: 2 minuti a domanda. Sembrano pochi, ma molte
domande si risolvono in 20 secondi (definizioni) e ti lasciano tempo per quelle lunghe
(output di un programma). La scheda del corso dà tre esempi di domande vere:
\q{Qual è l'output del seguente programma?}, \q{Scegliere la corretta definizione di
una funzione virtuale pura}, \q{Cosa intendiamo dire quando diciamo che un algoritmo X è
asintoticamente più efficiente di un altro algoritmo Y?}. Sono i tre tipi di domanda che
troverai: \textbf{output}, \textbf{definizione}, \textbf{complessità}.

\subsection*{La tecnica delle due passate}

\begin{passopasso}[: come gestire l'ora del quiz]
\begin{enumerate}
  \item \textbf{Prima passata (circa 25 minuti).} Leggi tutte le domande in ordine.
        Rispondi subito a quelle che sai con sicurezza. Se una domanda ti fa esitare
        più di un minuto, segnala con un cerchietto e vai avanti. Non restare bloccato:
        ogni domanda vale come le altre.
  \item \textbf{Seconda passata (circa 25 minuti).} Torna alle domande cerchiate. Ora
        hai tempo per simularle sul foglio: disegna la memoria, fai la tabella delle
        variabili.
  \item \textbf{Controllo (ultimi 10 minuti).} Verifica di aver segnato una risposta per
        ogni domanda (se gli errori non tolgono punti, \textbf{non lasciare nulla in
        bianco}) e di averla segnata nella casella giusta.
\end{enumerate}
\end{passopasso}

\subsection*{Un metodo per ogni tipo di domanda}

\begin{center}\small
\begin{tabular}{p{3.4cm}p{6.6cm}p{2.6cm}}
\toprule
\textbf{Tipo di domanda} & \textbf{Metodo} & \textbf{Dove ripassare} \\
\midrule
Output di un programma & Fai una \textbf{tabella} con una colonna per variabile e
una riga per istruzione. Non \q{immaginare} l'output: scrivilo riga per riga. & M2, M3, M4 \\
Puntatori e reference & Disegna stack e heap con le frecce, come nel Modulo 3. & M3 \\
Costruttori e distruttori & Costruzione: dalla base alla derivata. Distruzione:
l'ordine inverso. Oggetti locali distrutti in ordine inverso di creazione. & M4, M9 \\
Ereditarietà e \texttt{virtual} & Chiediti: il metodo è \texttt{virtual}? Si chiama
tramite puntatore/reference alla base? Solo se entrambe le risposte sono sì decide il tipo
\emph{vero} dell'oggetto. & M9 \\
Complessità & Conta quanti cicli annidati dipendono da $N$; se il problema si dimezza a
ogni passo è $\log N$. Usa la tabella del Modulo 6. & M6, M7 \\
Ordinamenti & Ricorda la tabella migliore/medio/peggiore e simula due passate a mano. & M7 \\
Ricorsione & Le tre domande del Modulo 8; scrivi le chiamate una sotto l'altra. & M8 \\
Strutture dati & Disegna i nodi. Pila: LIFO. Coda: FIFO. BST: sinistra minore, destra
maggiore. & M11--M13 \\
Visite di alberi e grafi & Pre (radice prima), in (radice in mezzo, nel BST esce
ordinata), post (radice dopo). BFS: coda, per livelli. DFS: ricorsione, in profondità. & M13, M14 \\
Definizioni & Elimina le risposte assurde, poi scegli la più \emph{precisa}. & Glossario (A.4) \\
\bottomrule
\end{tabular}
\end{center}

\begin{consiglio}[: come eliminare le risposte sbagliate]
Quando non sei sicuro, non cercare la risposta giusta: \textbf{elimina quelle
sbagliate}. Di solito due risposte su quattro sono chiaramente fuori tema. Tra le due
rimaste, diffida di quelle con parole assolute (\q{sempre}, \q{mai}, \q{solo},
\q{qualsiasi}): in informatica le regole hanno quasi sempre un'eccezione. Passare da 4
a 2 possibilità raddoppia le tue probabilità anche quando tiri a indovinare.
\end{consiglio}

\begin{trappola}[: le domande \q{a trabocchetto} più comuni]
\begin{itemize}
  \item Un programma che \textbf{non compila}: una delle risposte è quasi sempre
        \q{errore di compilazione}. Prima di simulare, controlla: accesso a un membro
        \texttt{private}? Istanza di una classe astratta? Metodo non \texttt{const} su un
        oggetto \texttt{const}?
  \item Divisione tra interi: \texttt{7/2} vale \texttt{3}, non \texttt{3.5}.
  \item \texttt{delete} su un array allocato con \texttt{new[]}: serve \texttt{delete[]}.
  \item Passaggio per valore: la funzione modifica una \emph{copia}, l'originale non
        cambia.
  \item Ordine delle visite: in-order su un albero \emph{qualsiasi} non è ordinata; lo è
        solo su un BST.
\end{itemize}
\end{trappola}

\subsection*{Come allenarsi}

\begin{itemize}
  \item I \textbf{quiz di tutti i 14 moduli} (circa 140 domande) e i 20 della simulazione:
        sempre cronometrati, 2 minuti a domanda.
  \item Il \textbf{quaderno degli errori}: per ogni domanda sbagliata scrivi in una riga
        \emph{perché} l'hai sbagliata (non sapevo, ho letto male, ho calcolato male).
        Se l'errore più frequente è \q{ho letto male}, il problema non è lo studio ma la
        fretta.
  \item Le domande del gioco \q{Chi vuol essere milionario} del tutorato 2024 sono
        domande di teoria in stile quiz. Le trovi nella cartella del tutorato 2024 dentro
        \texttt{Fonti/GitCatalogue/}, in \texttt{Esercizi/Soluzioni/}.
\end{itemize}

\begin{inparole}
Due passate: prima quelle facili, poi quelle lunghe. Per l'output fai la tabella, per i
puntatori fai il disegno, per il resto elimina le risposte sbagliate. L'obiettivo non è
30/30: è stare comodamente sopra la soglia (punta a 22--23 risposte giuste).
\end{inparole}

% ---------------------------------------------------------------------
\section{Prova 2 -- Il laboratorio: un programma che funziona}
% ---------------------------------------------------------------------

Il laboratorio è la prova più selettiva. La regola d'oro viene direttamente dalla scheda
del corso: all'orale accede solo chi consegna una soluzione \textbf{funzionante}. Questo
cambia tutto il modo di lavorare.

\begin{analogia}[: costruire una casa abitabile]
Due muratori hanno due ore per costruire una casetta. Il primo costruisce subito un
monolocale con tetto, porta e finestre, e poi aggiunge stanze una alla volta: in ogni
momento, se suona la sirena, ha una casa in cui si può vivere. Il secondo costruisce le
fondamenta di una villa di dieci stanze, poi tutti i muri, e al suono della sirena non ha
ancora il tetto: ha più mattoni del primo, ma nessuna casa. All'esame devi essere il
primo muratore. \textbf{Un programma piccolo che compila e funziona vale più di un
programma grande che non compila.}
\end{analogia}

\subsection*{Prima dell'esame: cosa devi avere in testa}

Le tracce UNICT hanno quasi sempre la stessa struttura (lo vedi in
\texttt{Fonti/GitCatalogue/}: Animale, Biblioteca, Frutto, Studente\_Borsista, Veicolo,
gli esami del 2022). Ci sono due famiglie:
\begin{enumerate}
  \item \textbf{Gerarchia di classi}: una base astratta con metodi \texttt{virtual} puri,
        due o tre classi derivate, il metodo \texttt{put} con \texttt{operator<<}, un array
        di puntatori alla base riempito a caso con \texttt{rand}, e poi delle operazioni
        sull'array (contare, filtrare per tipo con \texttt{typeid} o
        \texttt{dynamic\_cast}, calcolare medie, lanciare eccezioni). Modulo 9.
  \item \textbf{Struttura dati template}: lista, pila, coda o BST su nodi dinamici, con
        inserimento, stampa e un'operazione \q{originale} (ad esempio la \q{distanza
        successore} del 29 settembre 2022: il numero di archi tra il nodo con valore $X$
        e il nodo del suo successore). Moduli 11--13.
\end{enumerate}
Spesso le due famiglie si combinano: una struttura dati che contiene oggetti della
gerarchia.

\begin{passopasso}[: il kit da sapere a memoria]
Devi saper scrivere \textbf{senza guardare}, in meno di 20 minuti ciascuno:
\begin{enumerate}
  \item lo scheletro della gerarchia con \texttt{put} e \texttt{operator<<} (il
        \q{template esami di laboratorio} del docente, nel Modulo 9);
  \item il riempimento di un array di \texttt{Base*} con \texttt{rand} e la sua
        liberazione con \texttt{delete};
  \item una lista template con inserimento in testa, in coda e in ordine, stampa,
        ricerca, cancellazione e distruttore (Modulo 11);
  \item pila e coda template (Modulo 12);
  \item un BST template con inserimento, ricerca, minimo, successore, visite e altezza
        (Modulo 13);
  \item una classe eccezione personalizzata e il blocco \texttt{try}/\texttt{catch}
        (Modulo 9).
\end{enumerate}
Il test è semplice: foglio bianco, cronometro, compilatore. Se ci metti più di 20 minuti
o devi guardare, rifallo il giorno dopo.
\end{passopasso}

\begin{consiglio}[: allenati nelle condizioni dell'esame]
In laboratorio non avrai Internet, non avrai l'AI, forse non avrai il tuo editor
preferito. Allenati così:
\begin{itemize}
  \item \textbf{Stacca il Wi-Fi} durante le tracce d'esame.
  \item Compila \textbf{da terminale}, così sai farlo su qualunque macchina:
        \texttt{g++ -Wall -Wextra -o prog main.cpp} e poi \texttt{./prog}
        (su Windows \texttt{prog.exe}).
  \item Abituati a leggere la documentazione del C++ invece di cercare la soluzione
        già pronta: è l'unico aiuto che avrai.
  \item Se la consegna è con \texttt{git}, prova almeno una volta tutto il giro:
        \texttt{git clone}, \texttt{git add}, \texttt{git commit -m \q{...}},
        \texttt{git push}.
\end{itemize}
\end{consiglio}

\subsection*{Durante l'esame: il piano minuto per minuto}

Questo piano è pensato per 2 ore. Se la prova dura di più, allarga le fasi centrali.

\begin{center}\small
\begin{tabular}{lp{11cm}}
\toprule
\textbf{Minuti} & \textbf{Cosa fare} \\
\midrule
0--10 & \textbf{Leggi tutta la traccia due volte}, con la penna. Sottolinea i nomi delle
        classi, degli attributi e dei metodi richiesti, e ogni \q{deve}. Dividi la traccia
        in punti numerati: sono i tuoi step. \\
10--25 & \textbf{Scheletro che compila}: tutte le classi con gli attributi, costruttori,
        \texttt{put}/\texttt{operator<<}, un \texttt{main} che crea un oggetto e lo
        stampa. Compila. Da questo momento hai un programma funzionante. \\
25--95 & \textbf{Un punto della traccia alla volta}. Per ogni punto: scrivi, compila,
        prova nel \texttt{main} con una stampa che dice cosa stai provando. Solo quando
        funziona passi al punto successivo. Ordine: prima i punti che sai fare. \\
95--110 & \textbf{Pulizia}: \texttt{delete} di tutto ciò che hai allocato, distruttore
        della base \texttt{virtual}, \texttt{main} che mostra chiaramente ogni punto. \\
110--120 & \textbf{Consegna}: ricompila da zero, esegui un'ultima volta, consegna.
        Non fare modifiche dell'ultimo minuto: se rompi qualcosa non hai il tempo di
        sistemarlo. \\
\bottomrule
\end{tabular}
\end{center}

\begin{trappola}[: compilare troppo di rado]
Lo sbaglio più costoso in laboratorio è scrivere cento righe e poi compilare. Gli errori
si moltiplicano, si mescolano e il compilatore ne stampa cinquanta. Compila \textbf{ogni
10--15 righe}: un errore trovato subito costa un minuto, cinquanta errori insieme costano
mezz'ora e tanta ansia.
\end{trappola}

\subsection*{Quando qualcosa va storto}

\begin{center}\small
\begin{tabular}{p{3.6cm}p{9.6cm}}
\toprule
\textbf{Problema} & \textbf{Cosa fare} \\
\midrule
Tanti errori di compilazione & Leggi \textbf{solo il primo} errore e la riga che
indica: gli altri sono spesso conseguenze del primo. Correggi, ricompila. \\
Errore incomprensibile nei template & Cerca il nome della \emph{tua} funzione e il numero
di riga del \emph{tuo} file nel messaggio. Controlla di aver scritto
\texttt{template <typename T>} prima di ogni metodo definito fuori dalla classe, oppure
definisci tutti i metodi \emph{dentro} la classe: è più sicuro. \\
Errori del linker (\texttt{undefined reference}) & Un metodo è dichiarato ma non definito,
oppure hai messo un template in un file \texttt{.cpp} separato. Con i template tieni
tutto in un file. \\
Il programma si chiude da solo (segmentation fault) & Quasi sempre un puntatore
\texttt{nullptr} o già cancellato. Metti una stampa (\texttt{cout << \q{qui 1} << endl;})
prima e dopo il punto sospetto per capire dove si ferma. Controlla i casi limite:
struttura vuota, un solo elemento. \\
Il programma non finisce mai & Un ciclo che non avanza (ti sei dimenticato
\texttt{p = p->succ}?) o una ricorsione senza caso base. Interrompi con Ctrl+C. \\
Un punto della traccia non ti riesce & Lascialo da parte, \textbf{commenta} il codice
che non compila (con una riga che spiega cosa manca) e vai avanti. Torna dopo. Un
punto mancante è meglio di un programma che non compila. \\
\bottomrule
\end{tabular}
\end{center}

\begin{consiglio}[: prima di consegnare]
\begin{itemize}
  \item Compila senza errori? (Un programma che non compila non ti porta all'orale.)
  \item Il distruttore della base è \texttt{virtual}? I metodi ridefiniti sono
        \texttt{virtual} nella base?
  \item Ogni \texttt{new} ha il suo \texttt{delete} (\texttt{delete[]} per gli array)?
        Le strutture dati liberano i loro nodi nel distruttore?
  \item Le classi con memoria dinamica rispettano la Regola dei Tre (o vietano la copia
        con \texttt{= delete})?
  \item Hai gestito i casi limite: struttura vuota, un solo elemento, elemento non
        trovato?
  \item L'output del \texttt{main} mostra \emph{chiaramente} che ogni punto della traccia
        funziona, con delle scritte che dicono cosa si sta stampando? (La traccia del
        2022 lo chiedeva esplicitamente: \q{L'output del main deve mostrare chiaramente
        [\ldots] il funzionamento}.)
  \item Niente STL dove è vietata.
\end{itemize}
\end{consiglio}

\begin{inparole}
Arriva con il kit a memoria. In laboratorio: leggi la traccia due volte, costruisci
subito uno scheletro che compila, aggiungi un punto alla volta compilando spesso, e
consegna sempre un programma che compila, anche se incompleto.
\end{inparole}

% ---------------------------------------------------------------------
\section{Prova 3 -- L'orale: dove si decide il voto}
% ---------------------------------------------------------------------

L'orale comincia dal tuo laboratorio: si compila e si esegue il progetto, poi ti
chiedono delle \textbf{scelte implementative}. Poi si passa alla teoria di tutto il
programma. È una conversazione, non un'interrogazione a memoria: il docente vuole capire
se hai \emph{capito} quello che hai scritto.

\subsection*{Parte 1: il tuo codice}

Le domande tipiche sono:
\begin{itemize}
  \item Perché questo metodo è \texttt{virtual}? Cosa succederebbe senza?
  \item Perché il distruttore della base è \texttt{virtual}?
  \item Qual è la complessità di questo metodo? (Conta i cicli, come nel Modulo 6.)
  \item Cosa succede se la lista è vuota? Se l'elemento non c'è?
  \item Perché hai usato \texttt{dynamic\_cast} (o \texttt{typeid}) e non l'altro?
  \item Hai liberato la memoria? Dove?
  \item Come miglioreresti questo punto?
\end{itemize}

\begin{trappola}[: il codice che non sai spiegare]
Se all'esame scrivi un pezzo di codice che funziona ma che non sai spiegare (per esempio
copiato a memoria senza averlo capito), all'orale emerge in due domande. Quando ti alleni
sulle tracce, per ogni riga chiediti: \q{saprei dire perché c'è?}. Se la risposta è no,
quella riga è il tuo prossimo argomento di studio.
\end{trappola}

\subsection*{Parte 2: la teoria}

Ogni modulo ha da 3 a 5 domande d'orale con risposta modello (box viola): sono circa 70
in tutto. Prepararle \textbf{ad alta voce} è molto diverso dal leggerle.

\begin{passopasso}[: la risposta in tre tempi]
Per ogni domanda di teoria, rispondi in tre tempi:
\begin{enumerate}
  \item \textbf{Una frase di definizione}: \q{Un metodo virtuale puro è un metodo
        dichiarato con \texttt{= 0}, senza implementazione nella classe base.}
  \item \textbf{Un esempio o un disegno}: \q{Per esempio \texttt{Forma} ha
        \texttt{area()} virtuale pura, e \texttt{Cerchio} e \texttt{Quadrato} la
        implementano.} Se c'è un foglio, disegna.
  \item \textbf{Perché serve / cosa comporta}: \q{La classe diventa astratta: non si
        può istanziare, serve come interfaccia comune per il polimorfismo.}
\end{enumerate}
Tre frasi, trenta secondi: una risposta completa. Se il docente vuole sapere di più, te
lo chiede.
\end{passopasso}

\begin{consiglio}[: come allenarsi per l'orale]
\begin{itemize}
  \item \textbf{Registrati con il telefono} mentre rispondi a una domanda, poi
        riascoltati. Ti accorgi subito di frasi vaghe e di \q{ehm}.
  \item \textbf{Spiega a qualcuno} che non sa programmare (un familiare, un amico): se
        riesci a farti capire da lui, hai capito davvero. Le analogie della dispensa
        (scatoloni, indirizzi, fila al cinema) servono anche a questo.
  \item \textbf{Studia in coppia}: con un compagno di corso interrogatevi a vicenda
        sulle domande dei moduli. È uno degli allenamenti più efficaci.
  \item Allenati sulle \textbf{domande di collegamento}, quelle che separano il 24 dal
        28: \q{Perché l'inserimento in un BST può costare $O(N)$?}, \q{Che differenza c'è
        tra la visita per livelli di un albero e la BFS di un grafo?}, \q{Cosa hanno in
        comune la ricorsione e una pila?}
\end{itemize}
\end{consiglio}

\begin{trappola}[: non sai la risposta]
Succede a tutti. Non inventare e non restare in silenzio: dì \q{Non ne sono sicuro, provo
a ragionarci} e ragiona ad alta voce partendo da quello che sai. Un ragionamento
corretto, anche incompleto, vale molto più di una definizione recitata a memoria o di un
silenzio.
\end{trappola}

% ---------------------------------------------------------------------
\section{Il voto: cosa separa un 18 da un 28}
% ---------------------------------------------------------------------

La scheda del corso descrive le fasce di voto. Leggerle con attenzione ti dice dove
mettere lo sforzo.

\begin{center}\small
\begin{tabular}{lp{6cm}p{5.4cm}}
\toprule
\textbf{Voto} & \textbf{Cosa dice la scheda} & \textbf{In pratica} \\
\midrule
Non approvato & Non risponde ad almeno il 60\% delle domande né svolge gli esercizi. &
Quiz sotto soglia o laboratorio che non funziona. \\
18--23 & Padronanza minima, collegamenti modesti, risolve semplici esercizi. &
Laboratorio funzionante ma incompleto; all'orale sa le definizioni. \\
24--27 & Buona padronanza, buoni collegamenti, esercizi con pochi errori. &
Laboratorio quasi completo; all'orale spiega il proprio codice e collega gli
argomenti. \\
28--30 e lode & Padronanza completa, collegamenti con spirito critico, esercizi
completi e senza errori. &
Laboratorio completo e pulito (memoria liberata, casi limite); all'orale ragiona sulle
scelte e sulle alternative. \\
\bottomrule
\end{tabular}
\end{center}

\begin{consiglio}[: dove guadagnare punti]
\begin{itemize}
  \item Il salto da \textbf{18 a 24} si fa in laboratorio: completare tutti i punti
        della traccia. Si ottiene con le tracce d'esame fatte a cronometro.
  \item Il salto da \textbf{24 a 28} si fa all'orale: spiegare le scelte e collegare gli
        argomenti (complessità delle operazioni, perché una struttura e non un'altra).
  \item La parola che ritorna in tutte le fasce alte è \textbf{collegamenti}: non basta
        sapere i moduli uno per uno, bisogna vedere come si parlano.
\end{itemize}
\end{consiglio}

% ---------------------------------------------------------------------
\section{Tutti gli aiuti che hai a disposizione}
% ---------------------------------------------------------------------

Molti studenti preparano l'esame da soli, con un solo libro. Hai molte più risorse di
quelle che pensi, e quasi tutte sono gratuite.

\begin{center}\small
\begin{tabular}{p{3.6cm}p{9.6cm}}
\toprule
\textbf{Risorsa} & \textbf{Come usarla} \\
\midrule
Questa dispensa & Moduli, quiz, laboratori a step, 140 esercizi, simulazione d'esame,
formulario. È la base. \\
Tracce d'esame passate & In \texttt{Fonti/GitCatalogue/}: Animale, Biblioteca, Frutto,
Studente\_Borsista (anche nella versione con coda), Veicolo, gli esami del 2022 e
\q{Distanza successore}. Falle \textbf{a cronometro} senza guardare la soluzione. \\
Tutorato didattico & La scheda del corso prevede ore di lezione con un tutor
per \q{colmare eventuali lacune}. Partecipa e porta le tue domande. Nel 2024 il tutorato
ha fatto una lezione su \q{Consigli su come affrontare l'esame di laboratorio} e una
simulazione in aula. \\
Ricevimento del docente & Per le regole d'esame (vedi le tre domande all'inizio del
capitolo) e per i dubbi di teoria che non riesci a sciogliere. Vai con domande precise:
\q{non capisco perché qui serve \texttt{virtual}} è meglio di \q{non capisco
l'ereditarietà}. \\
Piattaforma Coding Contest & Consigliata dal docente per esercitarsi e autovalutarsi.
Alcuni esercizi sono in \texttt{Fonti/GitCatalogue/Programmazione2-main/.../Coding
Contest/}. \\
Simulazioni ed esercitazioni del corso & Il docente fa esercitazioni e simulazioni
d'esame durante il corso: se puoi, partecipa. \\
Compagni di corso e gruppi & Gruppi di studio e gruppi Telegram del corso: per sapere
com'è andato l'ultimo appello, che tipo di traccia è uscita, e per interrogarvi a
vicenda. \\
Libri di testo & Aguilar (testo 1) e Cormen (testo 2), citati nella scheda. Utili per un
argomento che non ti torna, non da leggere tutti. \\
Misure compensative & Se hai una certificazione DSA o di disabilità, contatta con
anticipo il docente e il referente CInAP del DMI: hai diritto a misure compensative
(per esempio più tempo). \\
\bottomrule
\end{tabular}
\end{center}

\begin{consiglio}[: andare a vedere un appello]
Se riesci, partecipa da \textbf{spettatore} agli orali di un appello prima del tuo (di
solito gli orali sono pubblici; chiedi conferma). In un'ora capisci che domande fa il
docente, quanto approfondisce e cosa apprezza. È l'informazione più preziosa che puoi
avere e nessuno la usa.
\end{consiglio}

% ---------------------------------------------------------------------
\section{Piano B: se una prova va male}
% ---------------------------------------------------------------------

Un esame che va male non è una sconfitta: è la simulazione più realistica che potessi
fare. Molti superano Programmazione 2 al secondo tentativo, e quasi sempre con un voto
migliore. L'importante è usare l'esperienza.

\begin{passopasso}[: subito dopo un appello andato male]
\begin{enumerate}
  \item \textbf{Il giorno stesso}, scrivi tutto quello che ricordi: le domande del quiz
        che ti hanno messo in difficoltà, la traccia del laboratorio, le domande
        dell'orale. Il giorno dopo l'avrai già dimenticato.
  \item \textbf{Capisci dove si è rotto}: quiz (teoria), laboratorio (velocità o
        mancanza di allenamento), orale (spiegare a voce)?
  \item \textbf{Rifai a casa la traccia del laboratorio} con calma, fino a farla
        funzionare. Poi rifalla a cronometro.
  \item \textbf{Controlla la data del prossimo appello} e la scadenza della prenotazione
        su SmartEdu, e prenotati subito.
  \item Ricalibra il piano di studio sulle debolezze che hai scoperto, non su tutto.
\end{enumerate}
\end{passopasso}

\begin{center}\small
\begin{tabular}{p{3.4cm}p{9.8cm}}
\toprule
\textbf{Cosa è andato storto} & \textbf{Cosa fare per l'appello successivo} \\
\midrule
Quiz non superato & Rifai tutti i quiz dei moduli e annota il \emph{tipo} di errore. Se
sbagli gli output, allenati con la tabella delle variabili; se sbagli le definizioni,
usa formulario e glossario. \\
Laboratorio non funzionante & Quasi sempre è un problema di \textbf{metodo}, non di
conoscenze: troppo codice scritto senza compilare, oppure un punto difficile affrontato
per primo. Fai 3--4 tracce in più seguendo il piano minuto per minuto. \\
Orale debole & Hai studiato leggendo e non parlando. Riprendi tutte le domande d'orale
\emph{ad alta voce}, registrandoti. \\
Ansia & Le simulazioni a cronometro servono proprio a questo: più la situazione è
familiare, meno fa paura. Fai la simulazione completa (quiz, laboratorio e orale) di
nuovo, nelle stesse condizioni. \\
\bottomrule
\end{tabular}
\end{center}

\begin{trappola}[: rinunciare a presentarsi]
Non saltare un appello perché \q{non sono ancora pronto al 100\%}: nessuno lo è mai. Se
hai fatto il piano di studio e almeno 5--6 tracce d'esame, presentati. Nel peggiore dei
casi hai fatto la migliore simulazione possibile e sai esattamente cosa ti aspetta.
\end{trappola}

% ---------------------------------------------------------------------
\section{Gli ultimi giorni e il giorno dell'esame}
% ---------------------------------------------------------------------

\subsection*{Gli ultimi tre giorni}

\begin{itemize}
  \item \textbf{Niente argomenti nuovi.} Ripassa formulario e glossario, rifai i quiz
        sbagliati, riscrivi a memoria il kit del laboratorio una volta.
  \item Controlla la \textbf{prenotazione} su SmartEdu, l'orario e l'aula. Guarda la
        pagina avvisi del docente: a volte cambiano aula o orario.
  \item Dormi. Il sonno è il momento in cui il cervello fissa quello che hai studiato:
        una notte in bianco prima dell'esame ti fa perdere più di quanto ti fa
        guadagnare.
\end{itemize}

\subsection*{La mattina dell'esame}

\begin{consiglio}[: cosa portare e cosa fare]
\begin{itemize}
  \item \textbf{Documento d'identità} (e tesserino universitario se lo hai).
  \item Penne (almeno due), acqua, qualcosa da mangiare: la giornata può essere lunga.
  \item Arriva con \textbf{almeno 15 minuti} di anticipo: arrivare di corsa ti fa
        cominciare il quiz con il cuore a mille.
  \item Al PC del laboratorio, prima di scrivere codice: apri un terminale, scrivi un
        \q{ciao mondo}, compilalo con \texttt{g++} ed eseguilo. In due minuti sai che
        l'ambiente funziona e sai come si usa.
  \item \textbf{Leggi tutta la traccia prima di scrivere}, e rileggila a metà prova.
\end{itemize}
\end{consiglio}

\begin{analogia}[: l'ansia è energia]
Prima di una gara anche gli atleti esperti hanno il cuore che batte forte: è il corpo
che si prepara a dare il massimo. Il problema non è l'ansia, è non sapere cosa fare. Se
hai un piano per ogni prova (quello di questo capitolo) e l'hai provato nelle
simulazioni, quando l'ansia arriva hai già una cosa concreta da fare: leggere la
traccia, costruire lo scheletro, compilare.
\end{analogia}

\clearpage
% ---------------------------------------------------------------------
\section{La checklist finale}
% ---------------------------------------------------------------------

Stampa questa pagina e spunta le voci man mano.

\begin{center}\small
\begin{tabular}{p{0.5cm}p{12.6cm}}
\toprule
\multicolumn{2}{l}{\textbf{Burocrazia}} \\
\midrule
$\square$ & Ho verificato la questione della frequenza obbligatoria (studente lavoratore). \\
$\square$ & Ho chiesto: penalità del quiz, documentazione consentita, modalità di consegna del laboratorio. \\
$\square$ & Sono prenotato su SmartEdu e conosco data, ora e aula. \\
\midrule
\multicolumn{2}{l}{\textbf{Quiz}} \\
\midrule
$\square$ & Ho fatto i quiz di tutti i 14 moduli a cronometro. \\
$\square$ & Ho fatto i 20 quiz della simulazione e supero il 60\% con margine. \\
$\square$ & So simulare un programma con la tabella delle variabili e disegnare stack e heap. \\
\midrule
\multicolumn{2}{l}{\textbf{Laboratorio}} \\
\midrule
$\square$ & So scrivere a memoria la gerarchia con \texttt{put}/\texttt{operator<<} e l'array di \texttt{Base*}. \\
$\square$ & So scrivere a memoria lista, pila, coda e BST template (20 minuti ciascuno). \\
$\square$ & Ho fatto almeno 5--6 tracce d'esame vere a cronometro, senza Internet. \\
$\square$ & So compilare ed eseguire da terminale con \texttt{g++} (e, se serve, consegnare con \texttt{git}). \\
\midrule
\multicolumn{2}{l}{\textbf{Orale}} \\
\midrule
$\square$ & Ho risposto ad alta voce a tutte le domande d'orale dei moduli. \\
$\square$ & So spiegare ogni riga delle mie soluzioni di laboratorio. \\
$\square$ & So rispondere alle domande di collegamento (complessità delle strutture, ricorsione e pila, BFS e visita per livelli). \\
\bottomrule
\end{tabular}
\end{center}

\begin{inparole}
Superare Programmazione 2 non richiede di essere un genio: richiede di conoscere le
regole, avere il kit del laboratorio a memoria, allenarsi nelle condizioni vere e
saper spiegare a voce quello che si è scritto. Se hai spuntato questa lista, sei pronto.
\end{inparole}
